\chapter{Analisi comparativa tra standard internazionali per l'adempimento di NIS2}
\label{cap:confronto}

% 3.1 %%%%%%%%%%%%%%%%%%%%%%%%%%%%%%%%%%%%%%%%%%%%%%%%%%%%%%%%%%%%%%%%%%%%%%%%
\section{Introduzione al confronto}
\subsection{ISO 27001, NIST CSF 2.0 e NIS2}
In questo capitolo dell'elaborato si svolge un confronto tra i framework presentati,
ISO 27001 e NIST CSF 2.0, e la Direttiva NIS2, con l'obiettivo di comprendere quale
ruolo essi possano avere nel soddisfare i requisiti stabiliti dal legislatore europeo.

L'analisi comparativa si basa sul confronto tra gli articoli NIS2 dal 20 al 25, le misure stabilite 
da ISO 27001 e le funzioni e i controlli del NIST CSF 2.0.
Il NIST CSF 2.0 e ISO 27001 presentano numerosi obiettivi comuni, ma possiedono differenze
significative nella struttura e nell'approccio adottati. Come evidenziato nella 
prassi di riferimento UNI/PdR 174, ISO 27001 ha un'impostazione sistemica ed è
strutturata a requisiti, in modo da consentire all'organizzazione la certificazione 
accreditata di un proprio SGSI. Il NIST CSF 2.0, invece, rappresenta uno strumento
per valutare e potenziare il proprio profilo di cybersicurezza in relazione al rischio
cui si è esposti. Di conseguenza, non tutti i controlli di ISO trovano una corrispondenza
in NIST CSF 2.0 e viceversa\footcite{uni-pdr174}. Tali divergenze emergono anche nel 
confronto con la Direttiva NIS2.


\subsection{Framework di cybersicurezza nazionali}
Un ruolo importante nell'adozione di NIS2 è ricoperto dai framework per la cybersicurezza varati
dagli Stati membri a livello nazionale. Le direttive europee, come visto, necessitano di un recepimento
nazionale per entrare in vigore e i singoli organi di governo in tale fase possono introdurre
ulteriori specifiche e modifiche alle misure stilate dall'Unione. In questo modo, si ottengono molteplici
derivazioni nazionali di NIS2, ognuna delle quali interpreta e rafforza la normativa in maniera leggermente
differente\footcite[p. 5]{parmar-miles-csf}, il che costringe le imprese multinazionali ad
adattare le proprie procedure per soddisfare i requisiti specifici per più nazioni\footcite[p. 18]{ecso-whitepaper-nis2}.

In questo contesto, è interessante notare come gli Stati membri incorporano gli standard internazionali
come NIST e ISO nel recepimento di NIS2, creando così framework nazionali ibridi.
Un approccio degno di nota è quello del Belgio, in cui la certificazione ISO 27001 è considerata
equivalente a soddisfare gli standard NIS2\footcite[p. 4]{ecso-whitepaper-nis2}. Altri Paesi, invece,
hanno sviluppato i propri framework nazionali a partire da quelli ISO e NIST. Tra essi vi è l'Italia,
che ha creato il \textit{Framework Nazionale per la Cybersecurity e la Data Protection} sulla 
base dell'adattamento del NIST CSF al contesto locale.

Il riconoscimento di framework esistenti come equivalenti a NIS2 rappresenta una
misura a grandissimo vantaggio delle aziende. Viste le numerose sfide nell'implementazione
di NIS2, esaminate nel Capitolo \ref{cap:nis2}, ciò porterebbe ad una considerevole riduzione
degli sforzi richiesti alle imprese, a vantaggio delle aziende più piccole e di chi 
già ha operato investimenti strategici in standard specifici\footcite[p. 8]{ecso-whitepaper-nis2}. 
Tuttavia è necessario riflettere sugli effetti a lungo e breve termine di una misura
di questo tipo. Infatti, se l'impatto sulle risorse aziendali nel breve termine è sicuramente
positivo, sul lungo periodo si individuano due potenziali rischi:
\begin{itemize}
    \item Lasciare agli Stati membri l'onere di stabilire quale framework riconoscere come
    valido porterebbe a un'ulteriore frammentazione delle metodologie di cybersicurezza aziendali,
    vanificando l'obiettivo di unificazione alla base di NIS2.
    \item Stabilire in maniera univoca un framework come equivalente a NIS2 porterebbe da un lato
    alla creazione di un monopolio di un unico ente certificatore e dall'altro all'accantonamento
    di NIS2, rendendola di fatto superflua.
\end{itemize}
Una strada percorribile potrebbe essere rappresentata dalla scelta di alcuni standard temporaneamente
equivalenti a NIS2, in modo da favorire le imprese che già li adottano e ridurre l'impatto
sulle risorse per l'implementazione delle misure di sicurezza, nell'ottica di un progressivo
adattamento ai requisiti NIS2. Ciò necessita della creazione di strumenti di mappatura dedicati
tra gli standard e NIS2\footcite[pp. 7--9]{ecso-whitepaper-nis2}, iniziativa a cui questo capitolo
ambisce a contribuire.

Un'altra soluzione di natura opposta è rappresentata dalla creazione di un unico framework di cybersicurezza 
europeo, proseguendo quindi nella direzione tracciata dal legislatore con NIS2, corredato da
framework per la gestione del rischio e della sicurezza della \textit{supply chain}, con l'obiettivo di
ridurre al minimo la frammentazione\footcite[pp. 4--6]{parmar-miles-csf}\footcite[pp. 7--9]{ecso-whitepaper-nis2}
e di accorpare in esso le disposizioni derivanti anche da altre normative rilevanti, quali GDPR e DORA.


\subsection{Tabella comparativa}
In questo capitolo il confronto è condotto con un’impostazione teorico–argomentativa: si chiariscono criteri, affinità e differenze tra NIS2, 
ISO/IEC 27001 e NIST CSF 2.0. A supporto e verifica di tale analisi, nel Capitolo \ref{cap:mapping} è riportata una tabella di mappatura che 
delinea l'allineamento delle misure prescritte dall'Art. 21 della Direttiva NIS2 con gli obiettivi esposti nell'allegato del framework NIST
e i controlli dell'\textit{Annex A} di ISO 27001.
La tabella traduce quanto discusso di seguito in evidenze implementative, consentendo di leggere sia il posizionamento dei singoli standard 
rispetto a NIS2 sia il potenziale effetto combinato nel percorso di attuazione.



% 3.1 %%%%%%%%%%%%%%%%%%%%%%%%%%%%%%%%%%%%%%%%%%%%%%%%%%%%%%%%%%%%%%%%%%%%%%%%
\section{Allineamento di ISO 27001 con NIS2}

\subsection{Convergenze}
La certificazione ISO 27001 è, come detto, tra le più adottate dalle aziende, con più di 40.000 certificazioni rilasciate a livello 
mondiale\footcite[p. 2]{mcintosh}. Essa garantisce evidenti vantaggi alle aziende rispetto agli 
adempimenti richiesti da NIS2, fornendo una solida struttura per la governance della cybersicurezza dell'azienda e l'adozione di 
controlli di sicurezza sui sistemi informatici\footcite[p. 35]{ecso-whitepaper-nis2}.

Un primo elemento di convergenza tra la Direttiva NIS2 e lo standard ISO/IEC 27001:2022 riguarda la centralità della governance e del ruolo del top management.
ISO impone alla direzione aziendale obblighi espliciti in materia di leadership, definizione delle policy e attribuzione di ruoli e responsabilità 
(clausole 5.1--5.3), che trovano un diretto parallelo in NIS2\footcite[artt. 20--21]{nis2-directive}, dove la responsabilità 
ultima della sicurezza informatica è attribuita agli organi societari.
Questo aspetto rende lo standard utile a strutturare la governance dell'azienda, fornendo un'organizzazione dettagliata, producendo documentazione
utile a certificare l'adozione di misure di cybersicurezza e facilitando l’accountability anche in ambito NIS2\footcite[pp. 56--58]{kouns-ciso}\footcite[p. 35]{ecso-whitepaper-nis2}, 
come sostenuto anche dai CISO intervistati.

Un secondo punto riguarda l'approccio alla gestione del rischio. ISO 27001 prevede un processo strutturato di valutazione e trattamento del rischio 
(clausole 6.1.2--6.1.3), ciclicamente riesaminato e aggiornato. Tale approccio è coerente con l’impostazione 
\textit{risk-based} di NIS2, che richiede analisi periodiche delle minacce e adeguamento delle misure di sicurezza\footcite[art. 21(2)]{nis2-directive}. 
La prassi UNI/PdR 174 sottolinea come l’adozione di ISO permetta di armonizzare le attività di \textit{risk assessment} con i requisiti operativi 
derivanti da NIS2\footcite[p. 9]{uni-pdr174}, costruendo una base per le attività di cybersicurezza fortemente improntata all'individuazione, 
alla valutazione e alla gestione del rischio.

Inoltre, i controlli operativi di ISO coprono altre aree chiave previste dalla direttiva. Le misure sulla gestione degli incidenti (A.5.24--A.5.28), 
il monitoraggio continuo (A.8.16) e le pratiche di \textit{threat intelligence} (A.5.07) trovano corrispondenza con le prescrizioni NIS2 sulla rilevazione, 
risposta e notifica degli incidenti. Anche le misure di continuità e ripristino (A.5.29--A.5.30, A.8.13--A.8.14) sono 
in linea con gli obblighi di \textit{business continuity} e \textit{disaster recovery} richiesti dalla direttiva\footcite[p. 5]{ecso-whitepaper-nis2}.

Le convergenze non si limitano agli aspetti tecnici, ma si estendono anche alla cultura organizzativa. ISO 27001 valorizza la formazione del personale (A.6.3) 
e la consapevolezza in materia di sicurezza, aspetti richiamati anche da NIS2 che introduce obblighi di istruzione e informazione,
oltre che di responsabilità.

Infine, i CISO intervistati spiegano come la motivazione principale per cui le aziende decidono di adottare ISO rimanga la certificabilità
delle misure. In questo modo, l'azienda dispone da un lato di un SGSI chiaro e ben strutturato e 
dall'altro di un valido strumento di accountability, ossia il certificato e tutti i documenti di monitoraggio
che ISO richiede. Questo permette di dimostrare, laddove un auditor esterno o il regolatore necessitino di verificare la conformità
dell'azienda ai requisiti, la bontà del sistema di gestione della sicurezza oggetto del controllo.


\subsection{Divergenze}
Pur essendo spesso reputata un punto di partenza essenziale per la compliance NIS2, il CISO di Prometeia precisa che 
la certificazione ISO 27001 non è completamente sovrapponibile alla direttiva.

In primo luogo, ISO non si applica ai sistemi dell'azienda in maniera integrata, ma il suo perimetro
è rappresentato da un singolo SGSI. L'azienda, spiega il CISO di Carrefour Italia, può decidere in seguito di certificare tutti i suoi SGSI oppure,
come spesso avviene, solamente quelli considerati più critici in relazione al profilo di rischio.
Al contrario, NIS2 interessa tutta l'azienda e i suoi sistemi di sicurezza per intero, essendo l'applicabilità
non volontaria, ma stabilita dal legislatore in base alla classificazione del soggetto come
\textit{essenziale} o \textit{importante}. Verosimilmente, dunque, molte aziende che adottano
ISO solo sui sistemi critici devono integrare misure ulteriori anche in aree che non dispongono di controlli informatici sufficienti.

Un tema importante della certificazione ISO, evidenziato dai CISO intervistati, è il fatto che ISO consenta di certificare misure di sicurezza, ma non 
garantisca l’efficacia sostanziale dei controlli in ottica NIS2.
ISO infatti adotta un approccio basato sul rischio, che lascia all’organizzazione la scelta dei controlli più adeguati in base al livello di rischio 
che l'azienda stessa ha ritenuto accettabile in relazione alle procedure di \textit{risk assessment} svolte. 
NIS2, al contrario, introduce istanze più prescrittive, rendendo di fatto obbligatorie alcune misure -- salvo motivata giustificazione. 
L'obiettivo di NIS2 è infatti quello di fornire alle imprese delle misure minime di resilienza,
ritenute indispensabili in relazione al livello di rischio dell'organizzazione all'interno 
del sistema Paese, richiedendo, come spiega Herin, una valutazione del rischio improntata su aspetti più concreti e non strettamente formali.
In questo senso, NIS2 ha un obiettivo più ampio e più operativo, imponendo la necessità di testare periodicamente
l'efficacia dei \textit{playbook} e delle misure di contrasto alle minacce, mentre ISO si concentra principalmente
sull'adozione formale dei controlli senza preoccuparsi dell'efficacia degli stessi in fase di audit.

Un’altra divergenza dovuta alla natura regolatoria di NIS2 riguarda gli obblighi di notifica. ISO 27001 prevede processi di gestione degli incidenti e di reporting interno (A.5.24--A.5.28), 
i quali possono fornire da supporto, ma non impone scadenze stringenti né modalità di comunicazione verso autorità esterne. 
NIS2 stabilisce invece tempistiche tramite KPI precisi per l'informazione delle autorità competenti in caso di incidenti significativi\footcite[art. 23]{nis2-directive}.
Questo comporta per le aziende certificate ISO la necessità di integrare playbook e workflow dedicati per garantire la compliance normativa, integrando
la documentazione e le procedure interne già implementate con ISO.

Si rileva anche una discostamento nella definizione delle responsabilità del top management. ISO si limita a richiedere leadership e impegni formali (clausole 5.1--5.3), 
mentre NIS2 introduce responsabilità dirette e, in alcuni casi, sanzioni personali a carico degli amministratori in caso di mancata adozione delle misure 
adeguate\footcite[art. 20]{nis2-directive}, richiedendo un coinvolgimento più attivo e sostanziale del management nella supervisione della cybersicurezza.

Infine, sia NIS2 che ISO/IEC 27001 pongono attenzione alla supply chain security, ma in modo differente. 
ISO fornisce strumenti per formalizzare controlli e verifiche lungo la catena del valore (A.5.19--A.5.20), mentre NIS2 rende questo obbligo 
cogente e sanzionabile, richiedendo ancora una volta controlli effettivi più stringenti, attribuendone la 
responsabilità a ogni soggetto NIS2\footcite[art. 21]{nis2-directive}.


\subsection{Ruolo nell'implementazione di NIS2}
Ai fini dell'adozione di NIS2, ISO 27001 si configura come uno strumento di governo manageriale, capace di introdurre responsabilità
e accountability nella gestione della sicurezza informatica, aiutando le aziende a:
\begin{itemize}
    \item creare un SGSI strutturato e manutenibile;
    \item disporre di obiettivi di controllo solidi;
    \item produrre documentazione per auditor esterni.
\end{itemize}
Lo standard abilita quindi le imprese a sviluppare capacità organizzative e processi di 
gestione del rischio verificabili, trovando numerose corrispondenze nella Direttiva NIS2, ma non coincide con la piena conformità normativa: 
per i CISO, rappresenta un abilitatore che necessita inevitabilmente di integrazioni da individuare tramite la conduzione di una 
\textit{gap analysis} tra le misure implementate dall'azienda e gli obblighi di legge.

Il vantaggio principale che introduce ISO è la disponibilità di un certificato rilasciato da un auditor ISO che verifica 
l'impegno da parte dell'azienda nell'adozione di misure a tutela della sicurezza informatica. Tuttavia, ciò non dimostra la completa
aderenza alla normativa: la certificazione ISO 27001 rappresenta un metodo per le aziende di dimostrare accountability laddove un auditor NIS2 
richieda di verificare le misure in essere, ma si dimostra carente nel fissare un livello di sicurezza effettiva coerente con le richieste di NIS2.
In particolare, spiegano i CISO interpellati, ISO 27001 certifica l'adozione di determinati strumenti di sicurezza e di documentazione da un punto 
di vista formale, mentre la vocazione di NIS2 è più operativa e richiede controlli sostanziali. 

Infine, è necessario riflettere sul diverso obiettivo di NIS2 e ISO 27001. La capacità di ISO di fornire una struttura ben definita per la
governance della cybersicurezza e l'implementazione di controlli sui sistemi di gestione della sicurezza informatica è un punto di forza 
nella transizione verso NIS2. Tuttavia, la direttiva europea include obblighi più stringenti, come il controllo della \textit{supply chain}
e gli obblighi di notifica in caso di incidenti. 

Michele Rivieri osserva inoltre la possibilità di integrare ISO 27001
con altri standard ISO, come il 27701 per la gestione delle informazioni sulla privacy e il 22301 per la gestione della
\textit{business continuity}, i quali contengono a loro volta elementi abilitanti alla compliance della direttiva.
Secondo ENISA\footcite{enisa-implementing-guidance}, invece, altri standard rilevanti sono:
\begin{itemize}
    \item ISO 22313:2020 e ISO 22301:2019 per misure di continuità operativa;
    \item ISO 22361:2022 per controlli inerenti a sicurezza, resilienza e gestione delle crisi;
    \item ISO 21500:2020 e ISO 21500:2021 per la gestione di progetti, programmi e portfolio.
\end{itemize}



% 3.2 %%%%%%%%%%%%%%%%%%%%%%%%%%%%%%%%%%%%%%%%%%%%%%%%%%%%%%%%%%%%%%%%%%%%%%%%
\section{Allineamento del NIST CSF 2.0 con NIS2}

\subsection{Convergenze}
Il NIST CSF 2.0 e NIS2 condividono una forte base contenutistica. Entrambi nascono, seppur in contesti differenti,
con l'obiettivo di fornire linee guida minime per la cybersicurezza e la resilienza informatica
e di contribuire allo sviluppo di un programma di cybersecurity nelle organizzazioni.

Una caratteristica comune tra NIST CSF e NIS2 è che entrambi richiedono artefatti attuativi. 
Il primo, infatti, è presentato come un insieme di \textit{best practices} che necessitano di una creazione
di un \textit{action plan} dedicato. La seconda, invece, deve essere codificata in un framework di cybersicurezza
a livello di ogni singolo Stato membro; sarà poi tale framework che le singole imprese utilizzeranno per derivare
il proprio, adattandolo alle esigenze specifiche del contesto\footcite{parmar-miles-csf}. Alternativamente, le imprese possono implementare
direttamente NIS2 nel proprio framework aziendale, ma con alcune implicazioni di cui si discute più avanti.

NIS2 e il NIST CSF riconoscono la governance come un costituente essenziale della strategia aziendale e della
gestione della sicurezza informatica nelle organizzazioni. La gestione dei processi, della tecnologia e delle 
risorse umane è infatti un elemento fondante delle misure prescritte dalle due fonti. All'interno della funzione
\textit{Govern} del NIST CSF trovano corrispondenza gli obblighi fissati da NIS2 per quanto riguarda\footcite[p. 4]{parmar-miles-csf}:
\begin{itemize}
    \item la creazione di una strategia di cybersicurezza uniforme, per favorire l'allineamento delle organizzazioni;
    \item la designazione di un punto di contatto all'interno delle organizzazioni, referente deputato alla gestione delle comunicazioni
    di cybersicurezza;
    \item la sicurezza della \textit{supply chain}, tramite la regolazione e il monitoraggio delle interazioni con entità parte 
    della catena del valore e la gestione degli incidenti che si possono verificare lungo essa.
\end{itemize}

Il ciclo delle funzioni del NIST CSF si svolge a partire dalla funzione \textit{Identify}. Qui trovano spazio le misure
relative alla gestione degli asset, l'identificazione delle vulnerabilità, il tracciamento delle informazioni di \textit{threat 
intelligence}, necessarie all'adempimento della Direttiva NIS2\footcite[p. 4]{parmar-miles-csf}.

Anche la funzione \textit{Protect} mostra un allineamento con NIS2, stabilendo controlli di sicurezza nelle aree di protezione
dei dati trasmessi, gestione delle identità digitali e specifici requisiti come l'autenticazione a due fattori e gli obblighi di
formazione delle risorse umane nell'ambito cyber\footcite[p. 4]{parmar-miles-csf}.

Per quanto riguarda le funzioni NIST inerenti alla protezione dalle minacce e al recupero da situazioni di crisi informatica,
il CSF e NIS2 continuano a dimostrarsi coerenti. \textit{Protect}, \textit{Respond} e \textit{Recover} forniscono misure che 
permettono all'organizzazione di strutturare una metodologia operativa per adempiere alle richieste del legislatore nei seguenti ambiti\footcite[p. 4]{parmar-miles-csf}:
\begin{itemize}
    \item predisposizione al dialogo con il CSIRT nazionale e le autorità in caso di incidenti significativi o di eventi di sicurezza informatica;
    \item necessità di sviluppare un \textit{Business Continuity Plan} (BCP) per codificare le procedure da adottare in caso di incidente 
    garantendo le funzionalità minime necessarie all'operatività dell'azienda;
    \item necessità di sviluppare un \textit{Disaster Recovery Plan} (DRP) per strutturare le operazioni per il ripristino della normale
    condizione di operatività in seguito a un incidente o un'interruzione grave dei servizi.
\end{itemize} 


\subsection{Divergenze}
Nonostante i numerosi punti di contatto, NIS2 e il NIST CSF presentano anche elementi di divergenza.

Il NIST CSF 2.0 è un framework creato dal \textit{National Institute of Standards and Technology} (NIST),
istituzione parte del governo degli Stati Uniti d'America. Per questo motivo, è strutturato sulla base del 
contesto americano, ma la sua struttura si adatta bene alle organizzazioni di tutto il mondo\footcite[pp. 1--3]{parmar-miles-csf}.
Al contrario, NIS2 è una direttiva dell'Unione Europea, pensata per l'applicazione nel contesto europeo. 

NIS2 pone infatti grande attenzione alla collaborazione a livello internazionale e alla condivisione di informazioni
riguardanti gli eventi di sicurezza informatica. L'idea alla base di tali misure è figlia della volontà del legislatore europeo di creare 
una community di esperti cyber per favorire l'innovazione e il progresso della sicurezza nelle organizzazioni, anche al di là dei confini 
nazionali. Ciò si presta particolarmente all'ecosistema europeo, individuando la necessità di ridurre la frammentazione e creare un fronte 
internazionale comune contro le minacce informatiche, vocazione che invece non si evince esplicitamente dal NIST CSF 2.0\footcite[p. 4]{parmar-miles-csf}.
Inoltre, data la sua natura regolatoria, NIS2 introduce un livello di specificità superiore per l'attuazione di determinate 
modalità operative rispetto a quanto avviene nel framework NIST. Si evidenzia infatti la presenza di precisi KPI 
(\textit{Key Performance Indicator}) fissati dal regolatore UE per la definizione di soglie precise per la gestione e la 
notifica degli incidenti alle autorità\footcite[p. 5]{parmar-miles-csf}, indicatori non previsti dal NIST.

Inoltre, le differenze strutturali riguardano anche la modalità con cui NIS2 e NIST CSF prescrivono le proprie misure.
Mentre il NIST CSF è un framework volontario e perciò l'adesione alle misure presentate è lasciata alla discrezione delle
aziende, NIS2 si configura come un insieme di obblighi a cui i soggetti interessati devono attenersi con rigore per dimostrare
il proprio adempimento e non incorrere in sanzioni (non previste nel caso del NIST CSF)\footcite[p. 4]{parmar-miles-csf}.

Un'ulteriore differenza fondamentale tra NIS2 e NIST CSF risiede nella modalità di implementazione. 
Seppur entrambi abbiano una vocazione fortemente operativa, superiore a ISO 27001, necessitano come visto di un ulteriore raffinamento per l'applicazione diretta.
Perciò il NIST suggerisce lo sviluppo di un piano d'azione che stabilisca i KPI (non presenti nel framework) e gli indicatori di rischio
utili alla valutazione delle misure da implementare sulla base delle vulnerabilità individuate. 
Al contrario, il processo di adozione di NIS2 può essere più complesso e rappresentare una sfida notevole, soprattutto per le 
organizzazioni con un livello di maturità dei sistemi di sicurezza inferiore. Per affrontare questo processo, 
può essere utile l'adozione di uno (o più) framework che forniscano una struttura solida su cui basare la strategia di sicurezza informatica\footcite{parmar-miles-csf}.

Infine, i CISO ascoltati concordano nell'individuare limite del framework NIST nella mancanza di una forma di certificazione riconosciuta delle misure implementate. 
Ciò riduce la rilevanza del CSF, rispetto a ISO 27001, nel dimostrare accountability per le misure di sicurezza informatiche, spingendo le aziende ad adottare
il NIST CSF solo parzialmente come strumento di supporto, data la grande coerenza con NIS2, ma senza la necessità di implementarlo nella sua interezza.


\subsection{Ruolo nell'implementazione di NIS2}
Viste le numerose analogie tra il framework NIST e NIS2 è evidente come il primo possa avere un ruolo fondamentale
nell'adozione della direttiva. Similmente a quanto visto con ISO, il NIST CSF permette alle organizzazioni di strutturare la governance,
i controlli e le procedure di sicurezza informatica, ottenendo una valida base da completare con gli adempimenti necessari
stabiliti da NIS2. In quest'ottica, può rappresentare una soluzione utile per aziende con strutture di sicurezza
informatica più acerbe, che necessitano di un approccio strutturato alle tematiche cyber, ma coadiuvato da integrazioni mirate.

Nella struttura del framework NIST Herin rileva una maggiore vocazione operativa, coerente con NIS2, rispetto allo standard ISO 27001. 
Nonostante il CSF sia presentato come ciclico, si osserva una divisione netta del framework in due parti:
\begin{enumerate}
    \item le funzioni \textit{Identify} e \textit{Protect} definiscono la condizione della sicurezza informatica durante
    la normale condizione di operatività, costruite sulla funzione \textit{Govern} che definisce la governance aziendale;
    \item le funzioni \textit{Detect}, \textit{Respond} e \textit{Recover} strutturano invece l'operatività aziendale
    all'avvenimento di un incidente di sicurezza e alla crisi che ne può derivare. Al raggiungimento dello step \textit{lessons learned}
    il ciclo si ripete a partire da \textit{Identify} per correggere il governo e le misure implementate sulla base delle analisi
    svolte sull'attacco subito e delle vulnerabilità individuate.
\end{enumerate}
Analogamente, NIS2 definisce sia le misure da implementare in condizioni di normale operatività sia quelle per la gestione
degli eventi di \textit{security breach}. Secondo il CISO di Carrefour, da questo punto di vista l'approccio NIST è più in linea 
con NIS2 rispetto a ISO 27001, definendo un insieme di processi utili a gestire operativamente tutti gli scenari cyber in azienda.

Tuttavia, è necessario integrare tali processi con le richieste dell'Unione per quanto 
riguarda gli obblighi di notifica e relazione con le autorità, non individuati dal NIST. 
In questo processo, è fondamentale tenere conto degli indicatori di rischio richiesti da NIS2 per identificare quali
misure siano da implementare in relazione al livello di rischio dell'azienda, nonché dei KPI utili in fase di gestione
delle emergenze di sicurezza informatica.
